\documentclass[10pt,a4paper]{amsart}
\usepackage{amsmath,amssymb,amsthm,graphicx}
% Retained preamble for the Stacks Project chapter snapshot in this directory.
% The chapter's own preamble.tex is a copy of the shared Stacks Project preamble; the block
% below is sliced MECHANICALLY, line for line, from it. Nothing was added, removed or
% redefined; the wrapper class, the xr-hyper/hyperref wiring and the externaldocument list
% are omitted because the assembler supplies the document class.
%
% lines taken from preamble.tex: 3, 5, 9, 11, 12, 156, 157, 158, ... (51 lines total)
%
% package coverage restored from the chapter's own preamble (2026-10-13): xy, xy-2cell

\usepackage[all]{xy}\usepackage{verbatim}
\xyoption{2cell}
\UseAllTwocells

\newenvironment{reference}{\comment}{\endcomment}
\newenvironment{slogan}{\comment}{\endcomment}
\newenvironment{history}{\comment}{\endcomment}
\theoremstyle{plain}
\newtheorem{theorem}[subsection]{Theorem}
\newtheorem{proposition}[subsection]{Proposition}
\newtheorem{lemma}[subsection]{Lemma}
\theoremstyle{definition}
\newtheorem{definition}[subsection]{Definition}
\newtheorem{example}[subsection]{Example}
\newtheorem{exercise}[subsection]{Exercise}
\newtheorem{situation}[subsection]{Situation}
\theoremstyle{remark}
\newtheorem{remark}[subsection]{Remark}
\newtheorem{remarks}[subsection]{Remarks}
\numberwithin{equation}{subsection}
\def\lim{\mathop{\mathrm{lim}}\nolimits}
\def\colim{\mathop{\mathrm{colim}}\nolimits}
\def\Spec{\mathop{\mathrm{Spec}}}
\def\Hom{\mathop{\mathrm{Hom}}\nolimits}
\def\Ext{\mathop{\mathrm{Ext}}\nolimits}
\def\SheafHom{\mathop{\mathcal{H}\!\mathit{om}}\nolimits}
\def\SheafExt{\mathop{\mathcal{E}\!\mathit{xt}}\nolimits}
\def\Sch{\mathit{Sch}}
\def\Mor{\mathop{\mathrm{Mor}}\nolimits}
\def\Ob{\mathop{\mathrm{Ob}}\nolimits}
\def\Sh{\mathop{\mathit{Sh}}\nolimits}
\def\NL{\mathop{N\!L}\nolimits}
\def\CH{\mathop{\mathrm{CH}}\nolimits}
\def\proetale{{pro\text{-}\acute{e}tale}}
\def\etale{{\acute{e}tale}}
\def\QCoh{\mathit{QCoh}}
\def\Ker{\mathop{\mathrm{Ker}}}
\def\Im{\mathop{\mathrm{Im}}}
\def\Coker{\mathop{\mathrm{Coker}}}
\def\Coim{\mathop{\mathrm{Coim}}}
\DeclareMathSymbol{\boxtimes}{\mathbin}{AMSa}{"02}
\def\QCohstack{\mathcal{QC}\!\mathit{oh}}
\def\Cohstack{\mathcal{C}\!\mathit{oh}}
\def\Spacesstack{\mathcal{S}\!\mathit{paces}}
\def\Quotfunctor{\mathrm{Quot}}
\def\Hilbfunctor{\mathrm{Hilb}}
\def\Curvesstack{\mathcal{C}\!\mathit{urves}}
\def\Polarizedstack{\mathcal{P}\!\mathit{olarized}}
\def\Complexesstack{\mathcal{C}\!\mathit{omplexes}}
\def\Pic{\mathop{\mathrm{Pic}}\nolimits}
\def\Picardstack{\mathcal{P}\!\mathit{ic}}
\def\Picardfunctor{\mathrm{Pic}}
\def\Deformationcategory{\mathcal{D}\!\mathit{ef}}
\usepackage{cleveref}
\crefname{originallemma}{original lemma}{original lemmas}
\Crefname{originallemma}{Original lemma}{Original lemmas}
\makeatletter
\newlabel{algebraic-lemma-map-space-into-stack}{{04T4}{source}{}{sourcepointer-1}{}}
\newlabel{algebraic-lemma-map-space-into-stack@cref}{{[originallemma][0][]04T4}{[1][0][]source}{}{}{}}
\makeatother
\crefname{originallemma}{original lemma}{original lemmas}
\Crefname{originallemma}{Original lemma}{Original lemmas}
\makeatletter
\newlabel{algebraic-remark-flat-fp-presentation}{{04WY}{source}{}{sourcepointer-2}{}}
\newlabel{algebraic-remark-flat-fp-presentation@cref}{{[originallemma][0][]04WY}{[1][0][]source}{}{}{}}
\makeatother
\crefname{originallemma}{original lemma}{original lemmas}
\Crefname{originallemma}{Original lemma}{Original lemmas}
\makeatletter
\newlabel{criteria-lemma-flat-quotient-flat-presentation}{{06FH}{source}{}{sourcepointer-3}{}}
\newlabel{criteria-lemma-flat-quotient-flat-presentation@cref}{{[originallemma][0][]06FH}{[1][0][]source}{}{}{}}
\makeatother
\crefname{originallemma}{original lemma}{original lemmas}
\Crefname{originallemma}{Original lemma}{Original lemmas}
\makeatletter
\newlabel{criteria-theorem-flat-groupoid-gives-algebraic-stack}{{06FI}{source}{}{sourcepointer-4}{}}
\newlabel{criteria-theorem-flat-groupoid-gives-algebraic-stack@cref}{{[originallemma][0][]06FI}{[1][0][]source}{}{}{}}
\makeatother
\crefname{originallemma}{original lemma}{original lemmas}
\Crefname{originallemma}{Original lemma}{Original lemmas}
\makeatletter
\newlabel{morphisms-lemma-finite-flat}{{02KB}{source}{}{sourcepointer-5}{}}
\newlabel{morphisms-lemma-finite-flat@cref}{{[originallemma][0][]02KB}{[1][0][]source}{}{}{}}
\makeatother
\crefname{originallemma}{original lemma}{original lemmas}
\Crefname{originallemma}{Original lemma}{Original lemmas}
\makeatletter
\newlabel{spaces-groupoids-lemma-quotient-stack-2-cartesian}{{04M9}{source}{}{sourcepointer-6}{}}
\newlabel{spaces-groupoids-lemma-quotient-stack-2-cartesian@cref}{{[originallemma][0][]04M9}{[1][0][]source}{}{}{}}
\makeatother
\crefname{originallemma}{original lemma}{original lemmas}
\Crefname{originallemma}{Original lemma}{Original lemmas}
\makeatletter
\newlabel{stacks-morphisms-lemma-properties-diagonal-from-presentation}{{0DTY}{source}{}{sourcepointer-7}{}}
\newlabel{stacks-morphisms-lemma-properties-diagonal-from-presentation@cref}{{[originallemma][0][]0DTY}{[1][0][]source}{}{}{}}
\makeatother
\crefname{originallemma}{original lemma}{original lemmas}
\Crefname{originallemma}{Original lemma}{Original lemmas}
\makeatletter
\newlabel{stacks-properties-lemma-check-property-covering}{{04XD}{source}{}{sourcepointer-8}{}}
\newlabel{stacks-properties-lemma-check-property-covering@cref}{{[originallemma][0][]04XD}{[1][0][]source}{}{}{}}
\makeatother
\crefname{originallemma}{original lemma}{original lemmas}
\Crefname{originallemma}{Original lemma}{Original lemmas}
\makeatletter
\newlabel{stacks-properties-lemma-quasi-compact-stack}{{04YC}{source}{}{sourcepointer-9}{}}
\newlabel{stacks-properties-lemma-quasi-compact-stack@cref}{{[originallemma][0][]04YC}{[1][0][]source}{}{}{}}
\makeatother
\title{Well-nigh affine iff a finite flat groupoid presentation exists}
\author{The Stacks Project Authors}
\date{stacks-project:stacks-more-morphisms.tex:local-snapshot-20261012; GNU FDL 1.2 or later}
\begin{document}
\maketitle
\noindent\textit{Editorial note.} Statement span is lines 3338-3348 of sources/stacks-stacks-more-morphisms/pinned-original-source.tex; proof span is lines 3350-3393, the complete author proof (42 body lines).\par\medskip
\noindent\textit{Editorial note.} Navigation stubs: 9 cited-but-not-included labels are supplied as non-load-bearing pointers with their pinned tags where the snapshot tag index has them.\par\medskip
\noindent\textit{Editorial note.} extra\_packages is empty on purpose: the retained preamble already loads xy with options and a second declaration clashes in the assembler wrapper.\par\medskip
\noindent\textit{Editorial note.} Text is retained verbatim, including the author's own typos and placeholders.\par\medskip
\noindent\textit{Editorial note.} Queue index 24 of queue\_stacks-stacks-more-morphisms.json (pb=2089, stmt 3338-3348, proof 3350-3393).\par\medskip
\noindent\textit{Editorial note.} stacks-stacks-more-morphisms wave smm01 (2026-10-14): fresh adjudication accept per lane-s3/smm\_judgments.md; pool per Chief 10:5x open-pool ruling (net 59 with-proof rows, bloodline lane-s3 by anchor audit, v2 true-key namespace stacks:stacks-more-morphisms.tex, row lemma-deform-property-fp-over-ft avoided on cross-chapter key collision); 60 with-proof rows judged 45 accept / 14 reject / 1 avoid; patch-sd-envs contract spans scanned.\par\medskip
\noindent\textit{Allowed standard prerequisites (editorial).}\begin{itemize}
\item Explicit prerequisite boundary. The cited inputs are prerequisites; the argument in the retained span is not replaced by any cited result.
\end{itemize}\medskip
\ifdefined\hypertarget\hypertarget{sourcepointer-1}{}\fi
\noindent\textit{Source reference.} Original lemma 04T4: sources/stacks-stacks-more-morphisms/pinned-original-source.tex, source lines 3393--3393. Cited by the retained statement or proof but not part of the collected spans; supplied only as a navigation pointer. The argument does not depend on its content.\par\medskip
\ifdefined\hypertarget\hypertarget{sourcepointer-2}{}\fi
\noindent\textit{Source reference.} Original lemma 04WY: sources/stacks-stacks-more-morphisms/pinned-original-source.tex, source lines 3393--3393. Cited by the retained statement or proof but not part of the collected spans; supplied only as a navigation pointer. The argument does not depend on its content.\par\medskip
\ifdefined\hypertarget\hypertarget{sourcepointer-3}{}\fi
\noindent\textit{Source reference.} Original lemma 06FH: sources/stacks-stacks-more-morphisms/pinned-original-source.tex, source lines 3393--3393. Cited by the retained statement or proof but not part of the collected spans; supplied only as a navigation pointer. The argument does not depend on its content.\par\medskip
\ifdefined\hypertarget\hypertarget{sourcepointer-4}{}\fi
\noindent\textit{Source reference.} Original lemma 06FI: sources/stacks-stacks-more-morphisms/pinned-original-source.tex, source lines 3393--3393. Cited by the retained statement or proof but not part of the collected spans; supplied only as a navigation pointer. The argument does not depend on its content.\par\medskip
\ifdefined\hypertarget\hypertarget{sourcepointer-5}{}\fi
\noindent\textit{Source reference.} Original lemma 02KB: sources/stacks-stacks-more-morphisms/pinned-original-source.tex, source lines 3393--3393. Cited by the retained statement or proof but not part of the collected spans; supplied only as a navigation pointer. The argument does not depend on its content.\par\medskip
\ifdefined\hypertarget\hypertarget{sourcepointer-6}{}\fi
\noindent\textit{Source reference.} Original lemma 04M9: sources/stacks-stacks-more-morphisms/pinned-original-source.tex, source lines 3393--3393. Cited by the retained statement or proof but not part of the collected spans; supplied only as a navigation pointer. The argument does not depend on its content.\par\medskip
\ifdefined\hypertarget\hypertarget{sourcepointer-7}{}\fi
\noindent\textit{Source reference.} Original lemma 0DTY: sources/stacks-stacks-more-morphisms/pinned-original-source.tex, source lines 3393--3393. Cited by the retained statement or proof but not part of the collected spans; supplied only as a navigation pointer. The argument does not depend on its content.\par\medskip
\ifdefined\hypertarget\hypertarget{sourcepointer-8}{}\fi
\noindent\textit{Source reference.} Original lemma 04XD: sources/stacks-stacks-more-morphisms/pinned-original-source.tex, source lines 3393--3393. Cited by the retained statement or proof but not part of the collected spans; supplied only as a navigation pointer. The argument does not depend on its content.\par\medskip
\ifdefined\hypertarget\hypertarget{sourcepointer-9}{}\fi
\noindent\textit{Source reference.} Original lemma 04YC: sources/stacks-stacks-more-morphisms/pinned-original-source.tex, source lines 3393--3393. Cited by the retained statement or proof but not part of the collected spans; supplied only as a navigation pointer. The argument does not depend on its content.\par\medskip
\begin{lemma}
\label{lemma-well-nigh-affine}
Let $\mathcal{X}$ be an algebraic stack. The following are equivalent
\begin{enumerate}
\item $\mathcal{X}$ is well-nigh affine, and
\item there exists a groupoid scheme $(U, R, s, t, c)$ with $U$ and
$R$ affine and $s, t : R \to U$ finite locally free such that
$\mathcal{X} = [U/R]$.
\end{enumerate}
If true then $\mathcal{X}$ is quasi-compact, quasi-DM, and separated.
\end{lemma}

\begin{proof}
Assume $\mathcal{X}$ is well-nigh affine. Choose an affine scheme $U$
and a surjective, flat, finite, and finitely presented morphism
$U \to \mathcal{X}$. Set $R = U \times_\mathcal{X} U$. Then we
obtain a groupoid $(U, R, s, t, c)$ in algebraic spaces and an
isomorphism $[U/R] \to \mathcal{X}$, see
Algebraic Stacks, Lemma \ref{algebraic-lemma-map-space-into-stack}
and Remark \ref{algebraic-remark-flat-fp-presentation}.
Since $s, t : R \to U$ are 
flat, finite, and finitely presented morphisms
(as base changes of $U \to \mathcal{X})$ we see that
$s, t$ are finite locally free
(Morphisms, Lemma \ref{morphisms-lemma-finite-flat}).
This implies that $R$ is affine (as finite morphisms are affine)
and hence (2) holds.

\medskip\noindent
Suppose that we have a groupoid scheme $(U, R, s, t, c)$ with $U$ and
$R$ are affine and $s, t : R \to U$ finite locally free.
Set $\mathcal{X} = [U/R]$. Then $\mathcal{X}$ is an algebraic stack
by Criteria for Representability, Theorem
\ref{criteria-theorem-flat-groupoid-gives-algebraic-stack} (strictly speaking
we don't need this here, but it can't be stressed enough that this is true).
The morphism $U \to \mathcal{X}$ is surjective, flat, and locally
of finite presentation by
Criteria for Representability, Lemma
\ref{criteria-lemma-flat-quotient-flat-presentation}.
Thus we can check whether $U \to \mathcal{X}$ is finite by
checking whether the projection $U \times_\mathcal{X} U \to U$
has this property, see Properties of Stacks, Lemma
\ref{stacks-properties-lemma-check-property-covering}.
Since $U \times_\mathcal{X} U = R$ by
Groupoids in Spaces, Lemma
\ref{spaces-groupoids-lemma-quotient-stack-2-cartesian}
we see that this is true. Thus $\mathcal{X}$ is well-nigh affine.

\medskip\noindent
Proof of the final statement. We see that $\mathcal{X}$
is quasi-compact by Properties of Stacks, Lemma
\ref{stacks-properties-lemma-quasi-compact-stack}.
We see that $\mathcal{X} = [U/R]$ is quasi-DM and separated by
Morphisms of Stacks, Lemma
\ref{stacks-morphisms-lemma-properties-diagonal-from-presentation}.
\end{proof}

\end{document}
